\section{A polynomially weighted Hurwitz resolvent}
\label{bhr:sec:weighted}

\subsection{The two Fourier directions and a convergent depth-two function}

For $z\in\C\setminus\R$, put
\begin{equation}\label{bhr:eq:coordinates}
 \begin{gathered}
 \epsilon=\sgn(\Im z),\quad \omega=2\epsilon\pi\ii,
 \quad \sigma=-\omega,\quad a=z+\epsilon\pi\ii,\\
 q=\frac{z-\epsilon\pi\ii}{z+\epsilon\pi\ii},\qquad
 c=\frac{2\epsilon\pi\ii}{a^2},\qquad |q|<1,\\
 \Log\omega=\log(2\pi)+\epsilon\pi\ii/2,\qquad
 \Log\sigma=\log(2\pi)-\epsilon\pi\ii/2.
 \end{gathered}
\end{equation}
Let $K(x)=\pi\cot(\pi x)$ and $R_z(x)=(K(x)-z)^{-1}$.
The elementary identity
\begin{equation}\label{bhr:eq:resolvent-fourier}
 R_z(x)=-\frac1a+c\sum_{n\ge1}q^{n-1}e^{\omega nx}
\end{equation}
converges uniformly with every argument derivative. In particular
$R_z(0)=R_z(1)=0$ and $R_z(x)=x+O(x^2)$ at zero.

The natural depth-two function for multiplication by a polynomial is
\begin{equation}\label{bhr:eq:omega-function}
 \mathcal O_{\alpha,\beta}(q)
   =\sum_{n,m\ge1}\frac{q^n}{(m+n)^\alpha m^\beta}
   =\Li_{\alpha,\beta}(q,q^{-1}).
\end{equation}
Here our convention is
$\Li_{\alpha,\beta}(u,v)=\sum_{k>m\ge1}
u^kv^m/(k^\alpha m^\beta)$.
The last expression in \eqref{bhr:eq:omega-function} is interpreted by
the displayed convergent series, not by assigning separate
independent branches to the two arguments.

\begin{lemma}[Normal convergence]\label{bhr:lem:omega-normal}
The series \eqref{bhr:eq:omega-function}, with all finite order
derivatives, converges normally on
\[
 |q|<1,\qquad \Re(\alpha+\beta)>1.
\]
It vanishes at $q=0$, and division by $q$ has a removable value there.
\end{lemma}
\begin{proof}
On a compact parameter set split the inner sum at $m=2n$.
For $m>2n$, the summand is bounded by a fixed multiple of
$|q|^n m^{-\Re(\alpha+\beta)}$. For $m\le2n$, all possible
growth is bounded by a fixed power of $n$, and the geometric
factor dominates it. Derivatives introduce only powers of
$\log m$, $\log(m+n)$ and $n$, which are handled by a small
margin in the strict inequalities. The same bounds justify
termwise differentiation. Every term contains $q^n$ with $n\ge1$.
\end{proof}

We shall also need the normally convergent subtracted functions
\begin{align}
 U_j(s,q)&=\sum_{n,m\ge1}q^n m^{-s}
       \bigl((m+n)^{-j}-m^{-j}\bigr),\label{bhr:eq:U}\\
 V_j(s,q)&=\sum_{n,m\ge1}q^n m^{-j}
       \bigl((m+n)^{-s}-m^{-s}\bigr),\label{bhr:eq:V}
\end{align}
where $j\ge1$. Each is holomorphic for $\Re s>-j$, $|q|<1$.
In the common initial domain they satisfy
\begin{equation}\label{bhr:eq:UV-continuation}
 U_j=\mathcal O_{j,s}-\frac{q}{1-q}\zeta(s+j),\qquad
 V_j=\mathcal O_{s,j}-\frac{q}{1-q}\zeta(s+j).
\end{equation}
For example, in the tail $m>2n$ each difference gains a factor
$O(n/m)$, uniformly on compact sets. The remaining part is
geometrically summable as in Lemma~\ref{bhr:lem:omega-normal}.
This proves the claimed larger domain directly, without assigning
values to divergent unsubtracted summands.

\subsection{The finite polynomial data}

For a polynomial $P$ of degree $d$, define
\begin{equation}\label{bhr:eq:polynomial-data}
 \mu_P=\int_0^1P(x)\,dx,\qquad
 b_j(P)=\frac{(-1)^{j-1}}{\omega^j}
    \bigl(P^{(j-1)}(1)-P^{(j-1)}(0)\bigr),\quad 1\le j\le d.
\end{equation}
The sum is empty when $P$ is constant. Repeated integration by
parts gives
\begin{equation}\label{bhr:eq:P-fourier}
 \int_0^1P(x)e^{\omega kx}\,dx
 =\begin{cases}
   \mu_P,& k=0,\\[2pt]
   \displaystyle\sum_{j=1}^d\frac{b_j(P)}{k^j},& k\in\Z\setminus\{0\}.
  \end{cases}
\end{equation}
The derivative of order $d$ is constant and hence contributes no
endpoint difference. Thus only these finitely many data occur.

Define the elementary polynomial resolvent
\begin{equation}\label{bhr:eq:Tpoly}
 T_P(z)=\int_0^1P(x)R_z(x)\,dx
       =-\frac{\mu_P}{a}+
          \frac cq\sum_{j=1}^d b_j(P)\Li_j(q).
\end{equation}
All quotients by $q$ in this article are given their removable
values at $q=0$, equivalently $z=\epsilon\pi\ii$.

\subsection{The master identity}

Set
\begin{align}\label{bhr:eq:Bpoly}
 \mathcal B_P(s,q)
   ={}&\sigma^{-s}\mu_P\Li_s(q)\notag\\
    &+\sum_{j=1}^d b_j(P)
       \left\{\omega^{-s}U_j(s,q)
        +\sigma^{-s}\bigl(\Li_s(q)\Li_j(q)
                              +(-1)^jV_j(s,q)\bigr)\right\}.
\end{align}

\begin{theorem}[Polynomial Hurwitz resolvent]\label{bhr:thm:weighted-master}
For every polynomial $P$, every $z\notin\R$, and $\Re s>0$,
\begin{equation}\label{bhr:eq:weighted-master}
 \boxed{\displaystyle
 I_P(s,z):=\int_0^1P(x)\zeta(1-s,x)R_z(x)\,dx
          =\frac{c\Gamma(s)}q\mathcal B_P(s,q).}
\end{equation}
The bracket $\mathcal B_P$ is holomorphic on $\Re s>-1$.
The formula therefore gives a meromorphic continuation through
$s=0$, with residue $-T_P(z)$ and no other pole in that half-plane.
On $\Re s>0$ it is a finite depth-two polylogarithmic formula
by \eqref{bhr:eq:UV-continuation}.
\end{theorem}

\begin{proof}
Initially take $\Re s>1$. The classical Hurwitz Fourier expansion
\cite[\S25.11]{bhr:DLMFZ} is absolutely convergent and reads
\begin{equation}\label{bhr:eq:hurwitz-fourier}
 \zeta(1-s,x)=\Gamma(s)\sum_{m\ge1}m^{-s}
       \left(\omega^{-s}e^{\omega mx}
                         +\sigma^{-s}e^{-\omega mx}\right).
\end{equation}
Multiply this by \eqref{bhr:eq:resolvent-fourier} and use
\eqref{bhr:eq:P-fourier}. The constant Fourier term of $R_z$ gives
\[
 -\frac{\Gamma(s)}a\sum_j b_j(P)
        \bigl(\omega^{-s}+(-1)^j\sigma^{-s}\bigr)\zeta(s+j).
\]
The equal-frequency diagonal contributes
$c\Gamma(s)\sigma^{-s}\mu_P\Li_s(q)/q$.
The same-sign frequencies contribute $\mathcal O_{j,s}$.
For the remaining frequencies split at $m=n$:
\begin{align}\label{bhr:eq:lattice-splitting}
 \sum_{\substack{n,m\ge1\\m\ne n}}
    \frac{q^n}{m^s(n-m)^j}
 &=\Li_s(q)\Li_j(q)+(-1)^j\mathcal O_{s,j}(q).
\end{align}
Indeed $m<n$ is parameterized by $n=m+h$ and gives the
product; $m>n$ is parameterized by $m=n+h$ and gives the second
sum. Absolute convergence justifies these changes of variables.

Insert \eqref{bhr:eq:UV-continuation}. Every ordinary zeta term
cancels because
\begin{equation}\label{bhr:eq:coefficient-cancellation}
 \frac{c}{1-q}=\frac1a.
\end{equation}
What remains is exactly \eqref{bhr:eq:weighted-master}.
Both sides are holomorphic on $\Re s>0$, so the initial identity
extends there. On the left, use
$\zeta(1-s,x)=x^{s-1}+\zeta(1-s,x+1)$:
the first term times $P R_z$ is locally integrable for $\Re s>-1$,
and the second has just the constant Hurwitz pole $-1/s$.
The normal convergence of $U_j,V_j$ proves the same assertion
on the right. Finally, telescoping in $m$ gives
\[
 U_j(0,q)=-\frac{\Li_j(q)}{1-q},\qquad V_j(0,q)=0.
\]
Since $\Li_0(q)=q/(1-q)$,
\begin{equation}\label{bhr:eq:Bzero}
 \mathcal B_P(0,q)=\mu_P\frac{q}{1-q}
                     -\sum_jb_j(P)\Li_j(q)=-\frac q c T_P(z).
\end{equation}
This proves the residue, including the removable $q=0$ case.
\end{proof}

\begin{example}[The first nonconstant weight]\label{bhr:ex:xweight}
For $P(x)=x$, one has $\mu_P=1/2$, $b_1=1/\omega$ and
\begin{align}
 \mathcal B_x(s,q)
  =\frac12\sigma^{-s}\Li_s(q)
   +\frac1\omega\left\{\omega^{-s}U_1(s,q)
       +\sigma^{-s}\bigl(\Li_s(q)\Li_1(q)-V_1(s,q)\bigr)\right\}.
 \label{bhr:eq:xweight}
\end{align}
This is the missing $x\log\Gamma(x)$ input to an individual
Barnes-$G$ transform. Setting $P=1$ deletes $U_j,V_j$ and
recovers the previously established unweighted formula.
\end{example}

\subsection{All weighted Stieltjes coefficients}

\begin{theorem}[A holomorphic Stieltjes generator]\label{bhr:thm:weighted-stieltjes}
For every polynomial $P$, integer $m\ge0$, and $z\notin\R$,
\begin{equation}\label{bhr:eq:weighted-stieltjes}
 \boxed{\displaystyle
 \int_0^1P(x)\gamma_m(x)R_z(x)\,dx
   =\frac{m!c}{q}[s^{m+1}]
               \Gamma(1+s)\mathcal B_P(s,q).}
\end{equation}
These are ordinary absolutely convergent integrals.
Equivalently their exponential generating function is
\begin{equation}\label{bhr:eq:weighted-stieltjes-germ}
 I_P(s,z)+\frac{T_P(z)}s
   =\frac cq\frac{\Gamma(1+s)\mathcal B_P(s,q)
                          -\mathcal B_P(0,q)}s.
\end{equation}
\end{theorem}
\begin{proof}
The shift law
$\gamma_m(x)=x^{-1}\log^m x+\gamma_m(x+1)$ and $R_z=O(x)$
prove convergence. The same bounds, locally uniformly for $s$
near zero, allow \eqref{bhr:eq:stieltjes-convention} to be integrated
after removing its constant pole. Now use Theorem~\ref{bhr:thm:weighted-master}
and \eqref{bhr:eq:Bzero}. The numerator on the right of
\eqref{bhr:eq:weighted-stieltjes-germ} vanishes at zero.
\end{proof}

For instance, the first two moments are obtained from
\begin{align}
 \int_0^1P\gamma_0R_z\,dx
  &=\frac cq\bigl(\partial_s\mathcal B_P(0,q)
                                      -\gamma\mathcal B_P(0,q)\bigr),\\
 \int_0^1P\gamma_1R_z\,dx
  &=\frac c{2q}\bigl(\partial_s^2\mathcal B_P(0,q)
       -2\gamma\partial_s\mathcal B_P(0,q)
                +(\gamma^2+\zeta(2))\mathcal B_P(0,q)\bigr).
\end{align}
Every derivative of the subtracted double sums is normally
convergent. In particular these formulas do not require subtracting
two divergent numerical evaluations of \eqref{bhr:eq:UV-continuation}.

\subsection{Argument derivatives and complete local corrections}

Let
\begin{equation}\label{bhr:eq:Ap}
 A_p(s)=\prod_{j=1}^p(s-j),\qquad
 r_{P,p}(z)=[x^p]\bigl(P(x)R_z(x)\bigr).
\end{equation}
The empty product is one. Meromorphic continuation of $I_P$ to
any desired finite spectral neighborhood can be constructed by
subtracting finitely many terms from the Taylor expansion of
$P(x)R_z(x)$ at zero. This supplies an unambiguous meaning to
the next formula even when $s-p$ is outside the normal domain
of the first subtraction in \eqref{bhr:eq:U}--\eqref{bhr:eq:V}.

\begin{theorem}[Weighted argument-derivative conversion]\label{bhr:thm:weighted-contact}
For $m\ge0$, $p\ge1$ and a polynomial $P$,
\begin{equation}\label{bhr:eq:weighted-contact}
 \boxed{\displaystyle
 \FP_x\int_0^1P(x)\gamma_m^{(p)}(x)R_z(x)\,dx
   =m![s^m]\left\{
       A_p(s)I_P(s-p,z)-\frac{r_{P,p}(z)A_p(s)}s\right\}.}
\end{equation}
The expression in braces is holomorphic near $s=0$.
The entire polynomial amplitude $A_p(s)$ must be subtracted,
not only its value at zero.
\end{theorem}
\begin{proof}
The Hurwitz argument ladder gives the meromorphic identity
\[
 \partial_x^p\zeta(1-s,x)=A_p(s)\zeta(1-s+p,x).
\]
The singular part at zero is $A_p(s)x^{s-p-1}$.
After multiplication by $P R_z$, the only exponent crossing
$-1$ at $s=0$ is $r_{P,p}A_p(s)x^{s-1}$.
Its continued integral is $r_{P,p}A_p(s)/s$, whereas the
finite part of every coefficient $x^{-1}\log^j x$ on $(0,1)$
is zero. All other terms commute with coefficient extraction.
To justify the latter statement, subtract a Taylor polynomial
through degree $p+1$ on $(0,\delta_0)$, integrate its terms
explicitly, and use normal integrability of the remainder.
The nonsingular Hurwitz tail at $x+1$ is analytic after an
argument derivative. This proves the formula and the
holomorphy assertion.
\end{proof}

For a completely explicit continuation of the depth-two expression,
one can subtract $N$ terms of $(1+n/m)^{-j}$ and $(1+n/m)^{-s}$.
The resulting known terms are respectively
\begin{align}
 &\sum_{k=0}^{N-1}\frac{(-1)^k(j)_k}{k!}
         \Li_{-k}(q)\zeta(s+j+k),\\
 &\sum_{k=0}^{N-1}\frac{(-1)^k(s)_k}{k!}
         \Li_{-k}(q)\zeta(s+j+k).
\end{align}
The remainders are normally convergent for $\Re s>1-j-N$.
This follows by the same $m>2n$ split and the finite Taylor
remainder estimate; the terms $m\le2n$ are still geometrically
summable. Thus \eqref{bhr:eq:weighted-contact} is a constructive
finite calculation plus convergent sums at every fixed $m,p$.

Repeated resolvents follow by $z$ differentiation:
\begin{equation}\label{bhr:eq:weighted-repeated}
 \FP_x\int_0^1\frac{P(x)\gamma_m^{(p)}(x)}{(K(x)-z)^r}\,dx
  =\frac1{(r-1)!}\partial_z^{r-1}
        \FP_x\int_0^1P(x)\gamma_m^{(p)}(x)R_z(x)\,dx.
\end{equation}
Local subtractions are holomorphic in $z$ on compact subsets of
either half-plane. If $P$ vanishes to order $e$ at zero, the
integral in \eqref{bhr:eq:weighted-repeated} is absolutely convergent
when $e+r>p$.
